% source: 464b908356409416ff1d3562cc4faf9bb61fc6e8
\chapter{中心二项分布（CBD）——为什么 Kyber 不使用高斯分布}
\chaptermark{中心二项分布（CBD）}
\label{ch:cbd}

\section{非常重要的一章}

在整个 Kyber 的故事中，本章是最容易被低估的章节之一。许多 Kyber 的入门介绍只说\textit{ CBD 不过是数几个比特而已}。这话没错，但它忽略了更深层的要点。真正的故事是：CBD 并不只是一个方便的技巧。它是一项设计选择，把格密码、实现安全、硬件效率和标准化历史联系在一起。

核心问题很简单：

\begin{quote}
为什么 NIST 最终为 ML-KEM 选择了 CBD，而不是更传统的离散高斯分布？
\end{quote}

读完本章，你应当能够解释：CBD 是什么，它为什么接近高斯分布，参数 \(\eta\) 的含义是什么，FIPS 203 的算法 2 如何运作，以及为什么 CBD 对常数时间实现如此有吸引力。

\section{回顾 LWE 的视角}

在 LWE 中，核心对象是噪声：

\[
\mathbf{b}=A\mathbf{s}+\mathbf{e}.
\]

安全性既不单独来自矩阵 \(A\)，也不单独来自秘密 \(\mathbf{s}\)。关键要素是噪声向量 \(\mathbf{e}\)。

换言之，整个密码学困难性都系于这样一个事实：误差项以一种难以恢复的方式隐藏了秘密。因此，为该误差选择一个好的采样分布至关重要。

\section{高斯采样的主要问题}

在最初的理论性 LWE 构造中，人们通常使用离散高斯分布。样本取自如下形式的分布：

\[
\Pr[X=x] \propto e^{-x^2/(2\sigma^2)}.
\]

这在数学上很优雅，但从工程角度看却很棘手。

实际问题在于，高斯采样通常需要代价高昂的操作：

\begin{itemize}
  \item 计算指数函数；
  \item 比较浮点数值；
  \item 拒绝采样；
  \item 反复循环，直到某个样本被接受。
\end{itemize}

这很慢，在软件中尤其如此。它也更难做到常数时间。

一个非常重要的副作用出现在实现安全方面。如果采样器使用拒绝循环和分支，运行时间就可能依赖于采样得到的值。这为计时攻击打开了大门：攻击者观察计时模式，试图推断出隐藏噪声的信息。

对于一个为现实部署而设计的密码系统来说，这是一个严重的缺陷。

\section{Kyber 的核心思想}

Kyber 并不试图以朴素的方式使用理想化的高斯分布。相反，它提出了一个更务实的问题：

\begin{quote}
能否构造一个看起来像高斯分布、却只由小整数、比特和简单算术构成的分布？
\end{quote}

答案是肯定的，这个答案就是 CBD。

CBD 简单、只使用整数，并且对常数时间实现非常友好。与高斯采样器相比，它在硬件和软件上都容易部署得多。

\section{最简单的例子}

我们从最小的非平凡情形开始：\(\eta=2\)。

假设我们采样四个随机比特：

\begin{verbatim}
1010
\end{verbatim}

把它们分成两组，每组两个比特：

\begin{verbatim}
10 10
\end{verbatim}

数出每一半中 1 的个数：

\begin{itemize}
  \item 前一半：\(10 \Rightarrow 1\)
  \item 后一半：\(10 \Rightarrow 1\)
\end{itemize}

然后相减：

\[
1-1 = 0.
\]

这就得到一个采样值。

再看另一个例子：

\begin{verbatim}
1110
\end{verbatim}

于是

\begin{itemize}
  \item 前一半：\(11 \Rightarrow 2\)
  \item 后一半：\(10 \Rightarrow 1\)
\end{itemize}

因此

\[
2-1 = 1.
\]

最后，取

\begin{verbatim}
0011
\end{verbatim}

于是

\begin{itemize}
  \item 前一半：\(00 \Rightarrow 0\)
  \item 后一半：\(11 \Rightarrow 2\)
\end{itemize}

所以

\[
0-2 = -2.
\]

由此，可能的取值为

\[
-2,-1,0,1,2.
\]

\section{为什么称为“中心”}

名称中的关键词是\textit{中心}。

在大量样本中，正的结果与负的结果相互抵消。例如，\(-2\) 和 \(2\) 这样的值很少出现，而 \(0\) 最为常见。该分布以零为中心。

事实上，均值恰好为零，因为两半在期望意义下的贡献相等。

所以该分布以零为中心，而这正是 LWE 噪声所需要的。

\section{为什么它越来越像高斯分布}

其优美的数学原因在于，每个比特都表现为一个 Bernoulli 随机变量：

\[
B_i \in \{0,1\}, \qquad \Pr(B_i=0)=\Pr(B_i=1)=\frac{1}{2}.
\]

现在考虑许多这样的 Bernoulli 变量之和。由中心极限定理，许多独立 Bernoulli 变量之和近似服从高斯分布。

CBD 正是建立在这一现象之上的。它不是一种粗糙的启发式方法。它是一个在数学上接近高斯分布的离散分布，因为它由两个二项计数之差构成。

所以真正的逻辑是：

\begin{quote}
CBD 并不只是在\textit{试图}看起来像高斯分布。它是一种离散构造，其形状自然地趋近于类高斯的钟形曲线。
\end{quote}

\section{\(\eta\) 的含义是什么？}

参数 \(\eta\) 控制样本的大小。

若 \(\eta=2\)，前一半使用两个比特，后一半也使用两个比特。

若 \(\eta=3\)，每一半各使用三个比特。

一般地，我们定义

\[
X = \sum_{i=1}^{\eta} B_i - \sum_{i=\eta+1}^{2\eta} B_i,
\]

其中每个 \(B_i\) 都是一个随机比特。

所得的取值范围为

\[
-\eta, -\eta+1, \dots, \eta.
\]

因此，\(\eta\) 越大，噪声幅度越大，分布也就越宽。

\section{ML-KEM 的实际参数选择}

在实际的 ML-KEM 参数集中，\(\eta_1\) 和 \(\eta_2\) 的取值并非随意。它们是根据安全性和解密失败率的要求来选取的。

\begin{table}[h]
\centering
\begin{tabular}{lcc}
\toprule
参数集 & \(\eta_1\) & \(\eta_2\) \\
\midrule
ML-KEM-512  & 3 & 2 \\
ML-KEM-768  & 2 & 2 \\
ML-KEM-1024 & 2 & 2 \\
\bottomrule
\end{tabular}
\caption{ML-KEM 中使用的典型 CBD 参数。}
\end{table}

这里 \(\eta_1\) 通常用于秘密的分布，而 \(\eta_2\) 用于加密噪声。

重要的是，这些选择并不是简单的\textit{越大越安全}。它们经过精心调校，以在安全性、正确性和实现效率之间取得平衡。

图~\ref{fig:cbd-pmf} 绘出了上表中实际出现的两个 $\eta$ 值所对应的精确概率质量函数，直接由 $\Pr[\mathrm{CBD}(\eta)=k]=\binom{2\eta}{\eta+k}/4^{\eta}$ 计算得到。

\begin{figure}[H]
\centering
\begin{subfigure}[t]{0.48\textwidth}
\centering
\begin{tikzpicture}
\begin{axis}[
  width=\textwidth, height=4cm,
  ybar, bar width=10pt,
  ymin=0, ymax=0.44,
  xtick=data,
  xlabel={$k$}, ylabel={$\Pr[X=k]$},
  symbolic x coords={-2,-1,0,1,2},
  label style={font=\scriptsize}, tick label style={font=\scriptsize},
  ymajorgrids=true, major grid style={pqcgray!25},
  axis lines*=left, enlarge x limits=0.15,
]
\addplot[fill=pqcblue, draw=pqcblue!70!black] coordinates {
  (-2,0.0625) (-1,0.25) (0,0.375) (1,0.25) (2,0.0625)
};
\end{axis}
\end{tikzpicture}
\caption{$\mathrm{CBD}(\eta=2)$}
\end{subfigure}
\hfill
\begin{subfigure}[t]{0.48\textwidth}
\centering
\begin{tikzpicture}
\begin{axis}[
  width=\textwidth, height=4cm,
  ybar, bar width=7pt,
  ymin=0, ymax=0.35,
  xtick=data,
  xlabel={$k$}, ylabel={$\Pr[X=k]$},
  symbolic x coords={-3,-2,-1,0,1,2,3},
  label style={font=\scriptsize}, tick label style={font=\scriptsize},
  ymajorgrids=true, major grid style={pqcgray!25},
  axis lines*=left, enlarge x limits=0.1,
]
\addplot[fill=pqcteal, draw=pqcteal!70!black] coordinates {
  (-3,0.015625) (-2,0.09375) (-1,0.234375) (0,0.3125) (1,0.234375) (2,0.09375) (3,0.015625)
};
\end{axis}
\end{tikzpicture}
\caption{$\mathrm{CBD}(\eta=3)$}
\end{subfigure}
\caption{ML-KEM-512/768/1024 中出现的两个 $\eta$ 值（见上表）所对应的中心二项分布的精确概率质量函数，由 $\Pr[\mathrm{CBD}(\eta)=k]=\binom{2\eta}{\eta+k}/4^{\eta}$ 计算得到：$\eta=2$ 时为 $\frac{1}{16},\frac{4}{16},\frac{6}{16},\frac{4}{16},\frac{1}{16}$，$\eta=3$ 时为 $\frac{1}{64},\frac{6}{64},\frac{15}{64},\frac{20}{64},\frac{15}{64},\frac{6}{64},\frac{1}{64}$。把 $\eta$ 从 $2$ 增大到 $3$ 会使钟形变宽变平——更多的概率质量从 $0$ 移开——代价是每个系数需要抽取两倍数量的随机比特。}
\label{fig:cbd-pmf}
\end{figure}

\section{FIPS 203 算法 2：发生了什么？}

在 FIPS 203 中，标准描述了一个采样过程，它本质上就是 CBD 采样操作。

其高层流程非常简单：

\begin{itemize}
  \item 输入一个长度为 \(2\eta\) 的比特串；
  \item 将其分成两组，每组长度为 \(\eta\)；
  \item 数出每组中 1 的个数；
  \item 将两个计数相减。
\end{itemize}

所以这个过程实际上就是：

\[
\text{CBD}_{\eta}(x) = \sum_{i=0}^{\eta-1} x_i - \sum_{i=\eta}^{2\eta-1} x_i.
\]

通俗地说：一个短的输入比特串被转换为一个小的有符号整数噪声样本。

这正是 Kyber 在实际实现中所依赖的操作。

如果用精确的比特下标、而不是上面“分成两组”这种非正式描述来精确地写出来，这就是 FIPS 203 的算法 2，即 \textsc{SamplePolyCBD}$_\eta$。它并不是每次从一个新的比特串中采样一个系数；而是消耗一个长度为 $64\eta$ 字节的长字节数组 $B$，并把它切分成 $256$ 个互不重叠的窗口，每个窗口 $2\eta$ 比特，每个输出系数对应一个窗口。

\begin{algorithm}[H]
\caption{SamplePolyCBD$_\eta$（FIPS 203，算法 2）}
\label{alg:samplepolycbd}
\begin{algorithmic}[1]
\Require 字节数组 $B \in \{0,\dots,255\}^{64\eta}$
\Ensure 多项式 $f \in \mathbb{Z}_q^{256}$，其系数为 $f_0,\dots,f_{255}$
\State $b \gets \textsc{BytesToBits}(B)$ \Comment{$b[0],\dots,b[8\cdot 64\eta-1]$；只用到前 $2\eta\cdot 256$ 个比特}
\For{$i = 0$ \textbf{to} $255$}
    \State $x \gets \sum_{j=0}^{\eta-1} b[2i\eta + j]$ \Comment{窗口 $i$ 的前一个 $\eta$ 比特半段中 1 的个数}
    \State $y \gets \sum_{j=0}^{\eta-1} b[2i\eta + \eta + j]$ \Comment{窗口 $i$ 的后一个 $\eta$ 比特半段中 1 的个数}
    \State $f_i \gets (x - y) \bmod q$
\EndFor
\State \Return $f = (f_0,\dots,f_{255})$
\end{algorithmic}
\end{algorithm}

因此，每个系数 $f_i$ 读取 $b$ 中属于它自己的一段 $2\eta$ 比特切片，起始比特偏移为 $2i\eta$：切片的前 $\eta$ 个比特计入 $x$，接下来的 $\eta$ 个比特计入 $y$，并且 $f_i=x-y \bmod q$。对于 $i=0$，这恰好就是上面手工演算过的“前一半、后一半、相减”的做法；唯一新增的内容是显式的偏移量 $2i\eta$，它对 $256$ 个系数中的每一个都精确指明了 $64\eta$ 字节数组中哪些比特属于它。

\section{一个简单的 Sage 实现}

下面是一个非常小的 Python 版本。

\begin{lstlisting}[language=Python]
from random import randint


def cbd(eta):
    a = 0
    b = 0

    for _ in range(eta):
        a += randint(0, 1)

    for _ in range(eta):
        b += randint(0, 1)

    return a - b
\end{lstlisting}

一个快速测试如下：

\begin{lstlisting}[language=Python]
for _ in range(20):
    print(cbd(2))
\end{lstlisting}

典型的输出可能是

\begin{verbatim}
0
1
-1
0
2
-2
\end{verbatim}

\section{直方图视角}

要真正理解 CBD，应当看一看它的直方图。

\begin{lstlisting}[language=Python]
samples = [cbd(2) for _ in range(100000)]
\end{lstlisting}

该分布趋向于集中在零附近，并且越来越呈钟形。这正是它常被描述为高斯分布的离散近似的原因。

用 \texttt{matplotlib} 可以做一个非常简单的可视化：

\begin{lstlisting}[language=Python]
from collections import Counter
import matplotlib.pyplot as plt

samples = [cbd(2) for _ in range(100000)]
cnt = Counter(samples)

plt.bar(cnt.keys(), cnt.values())
plt.show()
\end{lstlisting}

如果分别尝试 \(\eta=2\)、\(\eta=3\) 和 \(\eta=4\)，就可以看到形状变得更平滑、更对称。

\section{为什么标准不使用 \texttt{randint}}

一个微妙但重要的问题是，真正的 FIPS 实现并不会照字面使用 Python 中 \texttt{randint(0, 1)} 这类调用。

相反，标准通过诸如 SHAKE 这样的可扩展输出函数，从种子进行确定性扩展。

也就是说，种子被确定性地扩展：SHAKE 把它变成一个伪随机比特流，CBD 再把这些比特变成噪声样本。完整的流水线见第~\ref{sec:cbd-pipeline}~节。这一设计使噪声可以由种子复现，并且便于集成到更大的协议之中。

\section{为什么 CBD 对常数时间友好}

CBD 最大的实际优势之一，是它很容易以常数时间实现。

所涉及的操作只有：

\begin{itemize}
  \item 比特计数；
  \item 加法；
  \item 减法。
\end{itemize}

在基本流程中，既没有代价高昂的拒绝循环，也没有依赖数据的 \texttt{if} 分支。

这使得 CBD 比带有拒绝步骤的高斯采样器更容易防御计时攻击。

\section{对玩具版 Kyber 代码的一个小升级}

如果我们之前写过一个用简单例程生成小值的玩具版采样器，那么现在就可以把它替换为 CBD 式的采样。

例如，我们可以写成这样：

\begin{lstlisting}[language=Python]
def small_poly(n):
    coeff = []
    for _ in range(n):
        coeff.append(cbd(2))
    return coeff
\end{lstlisting}

这样一来，玩具版 Kyber 代码就更接近实际的 ML-KEM 设计了。

\section{FIPS 203 中完整的噪声流水线}
\label{sec:cbd-pipeline}

现在我们可以把 FIPS 203 中噪声所经历的完整路径拼接起来：从初始种子一直到一个模 LWE 样本。

\begin{figure}[H]
\centering
\begin{tikzpicture}[node distance=6mm and 6mm, every node/.style={font=\scriptsize}]
  \node[flownode] (seed) {seed};
  \node[flownode, right=of seed] (shake) {SHAKE / XOF};
  \node[flownode, right=of shake] (bits) {random\\[-2pt]bits};
  \node[flownode, right=of bits] (cbd) {CBD};
  \node[flownode, right=of cbd] (noise) {small-integer\\[-2pt]noise};
  \node[keynode, right=of noise] (mlwe) {Module-LWE};
  \draw[flowarrow] (seed) -- (shake);
  \draw[flowarrow] (shake) -- (bits);
  \draw[flowarrow] (bits) -- (cbd);
  \draw[flowarrow] (cbd) -- (noise);
  \draw[flowarrow] (noise) -- (mlwe);
\end{tikzpicture}
\caption{FIPS 203 中完整的噪声流水线。一个短种子经 SHAKE 扩展为伪随机比特流；CBD 把这些比特转换为小的有符号整数；这些整数成为模 LWE 样本的秘密和误差。每一步都是确定性的、常数时间的。}
\label{fig:cbd-pipeline}
\end{figure}

因此，噪声并不是什么神秘的浮点对象。它是一个小整数，通过安全的扩展过程由种子确定性地导出，并且沿途的每一步都是可复现的、常数时间的——这正是一个标准所要求的。

\section{小结}

本章的重要教训是：CBD 不只是一个方便的技巧。它是一种非常实用的采样方法，兼具四个有用的性质：

\begin{enumerate}
  \item 它的形状近似于高斯分布。
  \item 它易于实现。
  \item 它容易做到常数时间。
  \item 它非常适合软件和硬件部署。
\end{enumerate}

这就是为什么 CBD 成为 ML-KEM 中的标准选择，而不是理论上更为传统的高斯采样器。

一句话概括：

\begin{quote}
CBD 提供了一个对基于格的安全性而言足够好的噪声分布，它易于实现，而且比高斯采样更容易防御侧信道攻击。
\end{quote}

\section{下一步建议}

接下来自然的话题是内部的哈希与扩展层：为什么 SHAKE 和 Keccak 如此强大，以及一个短种子如何生成大量伪随机材料。这是从 CBD 通往实际 FIPS 203 和 FIPS 204 实现的桥梁。

\printbibliography[heading=chapterrefs]
